\section*{अध्याय \ref{ch:b1:nucleic-acids} --- न्यूक्लियोटाइड और न्यूक्लिक अम्ल}

\begin{solution}{exo:b1:nucleic-acids:1}
कोई \omterm{def:b1:nucleic-acids:nucleotide}{नाइट्रोजनी क्षारक}, कोई पेंटोस, एक से तीन फ़ॉस्फ़ेट। \omterm{def:b1:nucleic-acids:chain}{RNA}: राइबोस
(एक $2'$ हाइड्रॉक्सिल के साथ) और यूरैसिल; \omterm{def:b1:nucleic-acids:chain}{DNA}: $2'$-डिऑक्सीराइबोस और थाइमीन।
\end{solution}

\begin{solution}{exo:b1:nucleic-acids:2}
$5'$-TAGGATCC-$3'$। युग्म: G--C, G--C, A--T, T--A, C--G, C--G, T--A,
A--T: चार G--C (12 बंध) और चार A--T (8): यानी 20 \omterm{def:b1:water-small-molecules:water}{हाइड्रोजन बंध}।
\end{solution}

\begin{solution}{exo:b1:nucleic-acids:3}
T \qty{22}{\%}; G और C बची \qty{56}{\%} बाँटते हैं: हर एक \qty{28}{\%};
$G + C = \qty{56}{\%}$।
\end{solution}

\begin{solution}{exo:b1:nucleic-acids:4}
दोनों वक्रों के बीच, प्रति प्रतिशत \qty{0.4}{\celsius} पर: लगभग
\qty{85}{\celsius}। अवशोषकता \qty{40}{\%} चढ़ती है।
\end{solution}

\begin{solution}{exo:b1:nucleic-acids:5}
$\num{4.6e6}\times 0.34\,\mathrm{nm} = \qty{1.56}{mm}$; \num{460000}
घुमाव; लंबाई में $1560/2 = 780$ गुना (गुणसूत्र फंदों और अधिकुंडलों में
मुड़ा होता है)।
\end{solution}

\begin{solution}{exo:b1:nucleic-acids:6}
क्षार अपने किनारे आमने-सामने करके युग्मित होते हैं, और हर एक की शर्करा
युग्म के एक ही ओर तभी जुड़ती है जब दोनों मेरुदंड उलटी दिशाओं में चलें;
समांतर लड़ियों के साथ दोनों शर्कराएँ उलटे किनारों पर बैठतीं और युग्म की
चौड़ाई तथा मेरुदंड की स्थिति अनुक्रम के साथ बदलती रहती।
प्यूरीन--पिरिमिडीन युग्मों वाली प्रतिसमांतर लड़ियाँ वही अचर \qty{2}{nm}
देती हैं जो तंतु-प्रतिरूप माँगता है।
\end{solution}

\begin{solution}{exo:b1:nucleic-acids:7}
$0.35\times 50 = \qty{17.5}{\micro g/mL}$। $A_{260}/A_{280} = 1.75$:
यानी अनिवार्यतः शुद्ध \omterm{def:b1:nucleic-acids:chain}{DNA} (1.8), शायद प्रोटीन के एक अंश के साथ।
\qty{1}{mL} में: $17.5\times 10^{-6}/(650\times 1.66\times 10^{-24}) =
\num{1.6e16}$ \omterm{thm:b1:nucleic-acids:helix}{क्षारक युग्म}।
\end{solution}

\begin{solution}{exo:b1:nucleic-acids:8}
चट्टे में लगे \omterm{def:b1:nucleic-acids:nucleotide}{क्षारक} एक-दूसरे को ढाल देते हैं और कम अवशोषित करते हैं;
लड़ियाँ अलग करने से हर \omterm{def:b1:nucleic-acids:nucleotide}{क्षारक} प्रकाश के सामने खुल जाता है और अवशोषकता
\qty{40}{\%} चढ़ जाती है। कोई शुद्ध \omterm{def:b1:nucleic-acids:chain}{DNA} अपने एकमात्र $T_m$ से कुछ ही
डिग्री के भीतर सहकारी ढंग से पिघलता है; किसी मिश्रण में कई $G + C$
मात्राओं के \omterm{def:b1:nucleic-acids:chain}{DNA} होते हैं, हर एक का अपना $T_m$, इसलिए उनके वक्रों का योग
दसियों डिग्री भर फैल जाता है।
\end{solution}

\begin{solution}{exo:b1:nucleic-acids:9}
\qty{62}{\celsius} से ठीक नीचे और \qty{57}{\celsius} से ऊपर, मान लीजिए
\qty{60}{\celsius}: पूरा संकर टिका रहता है, बेमेल वाला
पिघल जाता है। \qty{45}{\celsius} पर तीन तक बेमेल वाले संकर
स्थायी रहते हैं और जाँच संबंधित अनुक्रमों को भी जगा देती है।
\end{solution}

\begin{solution}{exo:b1:nucleic-acids:10}
ये नियम दो पूरक लड़ियों के बीच क्षारक-युग्मन से निकलते हैं।
\omterm{def:b1:nucleic-acids:chain}{RNA}, जो एकलड़ी है, अपनी संरचना बाँधने वाला कोई साथी नहीं रखता, इसलिए आम
तौर पर $A \neq U$ होता है; और प्रश्न वाले छोटे विषाणुओं के जीनोम
एकलड़ी \omterm{def:b1:nucleic-acids:chain}{DNA} के हैं, जिन्हें भी इन नियमों को मानना ज़रूरी नहीं ---
और उनका इन्हें न मानना ही पहले संकेतों में था कि वे
एकलड़ी हैं।
\end{solution}

\begin{solution}{exo:b1:nucleic-acids:11}
अकेला हर युग्म प्रति मोल कुछ ही किलोजूल का है, जो तापीय
हलचल से कम है, और माइक्रोसेकंडों में बनता तथा टूटता है। किसी कुंडली में
युग्म साथ बनते हैं: हर युग्म पिछले पर चट्टा बनाता है और पहले युग्म बाकी
को जमा देते हैं, इसलिए ऊर्जाएँ जुड़ जाती हैं और सबके एक साथ टूटने की
प्रायिकता नगण्य हो जाती है --- यही सहकारिता है। चट्टा-बनना उतना ही
योगदान देता है जितना \omterm{def:b1:water-small-molecules:water}{हाइड्रोजन बंध}। किसी जाँच के लिए इसका अर्थ यह है कि
स्थायित्व लंबाई के साथ बढ़ता है और हर बेमेल के साथ तीखा गिरता है, जो उस
सहकारी दौड़ को तोड़ देता है: विशिष्टता समूचे खंड का गुण है।
\end{solution}

\begin{solution}{exo:b1:nucleic-acids:12}
कुंडली नकल की व्याख्या तुरंत कर देती है: हर लड़ी दूसरी को तय करती है,
इसलिए उन्हें अलग करके हर एक पर नए \omterm{def:b1:nucleic-acids:nucleotide}{न्यूक्लियोटाइड} युग्मित करने से दो
समरूप अणु मिल जाते हैं; और कोई उत्परिवर्तन एक नकल में जमा कोई गलत क्षार
है जिसकी उसके बाद वफ़ादारी से नकल होती रहती है। वह स्वयं यह नहीं बताती कि
चार क्षारों का कोई अनुक्रम किसी प्रोटीन को कैसे तय करता है --- उसमें एक
दशक लगा (कूट, mRNA, tRNA, राइबोसोम)। यह प्रतिरूप इसलिए स्वीकारा गया कि वह
हर माप से (शारगाफ़, X-किरण की मापें, घनत्व) बिना कुछ बचे बैठ गया, और
इसलिए भी कि आनुवंशिकता की उसकी व्याख्या इतनी मितव्ययी थी कि विकल्प ---
कोई संयोग --- अविश्वसनीय था; फिर मेसेल्सन और स्टाल ने उस नकल की पुष्टि
की जिसकी उसने भविष्यवाणी की थी।
\end{solution}

\begin{solution}{pb:b1:nucleic-acids:1}
\textbf{1.} $\num{48502}\times 0.34 = \qty{16490}{nm} = \qty{16.5}{\micro m}$।
\textbf{2.} $\num{48502}/10 = 4850$ घुमाव।
\textbf{3.} $\num{48502}\times 650 = \qty{3.15e7}{Da}$।
\textbf{4.} $\num{3.15e7}\times 1.66\times 10^{-24} = \qty{5.2e-17}{g}$।
\textbf{5.} प्रति युग्म दो: \num{97004} फ़ॉस्फ़ेट, शुद्ध आवेश $-97\,004\,e$।
\textbf{6.} $G + C = 0.499\times\num{97004} = \num{48405}$: $G = C =
\num{24203}$; $A = T = \num{24299}$।
\textbf{7.} $3\times\num{24203} + 2\times\num{24299} = \num{121207}$।
\textbf{8.} $81.5 + 16.6\times(-1) + 0.41\times 49.9 - 500/48502 =
81.5 - 16.6 + 20.5 - 0.01 = \qty{85.4}{\celsius}$।
\textbf{9.} $81.5 - 33.2 + 20.5 = \qty{68.8}{\celsius}$। धनायन दोनों
ऋणावेशित मेरुदंडों के बीच का प्रतिकर्षण ढाल देते हैं; कम
लवण के साथ लड़ियाँ अधिक प्रतिकर्षित करती हैं और नीचे तापमान पर अलग हो जाती हैं।
\textbf{10.} \qty{35}{\%}: $81.5 - 16.6 + 14.4 - 0.5 =
\qty{78.8}{\celsius}$; \qty{62}{\%}: $81.5 - 16.6 + 25.4 - 0.5 =
\qty{89.8}{\celsius}$। समूचा अणु चरणों में पिघलता है: A--T-धनी
क्षेत्र पहले खुलते हैं, G--C-धनी सबसे बाद में, इसलिए वक्र किसी एक तीखे
चरण के बजाय कंधों वाली एक चौड़ी चढ़ाई है।
\textbf{11.} $0.50\times 50 = \qty{25}{\micro g/mL}$; पिघलने के बाद
$A_{260} = 0.70$।
\textbf{12.} $25\times 10^{-6}/5.2\times 10^{-17} = \num{4.8e11}$ अणु।
\textbf{13.} $81.5 - 16.6 + 0.41\times 50 - 500/12 = 81.5 - 16.6 + 20.5
- 41.7 = \qty{43.7}{\celsius}$ (इतने छोटे खंडों के लिए यह सूत्र मोटा
है)। \qty{37}{\celsius} पर सिरे युग्मित तो होते हैं पर खुलकर साँस लेते;
\omterm{def:b1:cell-unit-of-life:cell}{कोशिका} में कोई एंजाइम (लाइगेज़) दोनों खाँचों को सहसंयोजी ढंग से सील कर
देता है, और तब वह वलय उन बारह युग्मों पर निर्भर नहीं रहता।
\textbf{14.} $\pi\times(1.0)^2\times\num{16490} = \qty{5.2e4}{nm^3}$।
\textbf{15.} $\frac{4}{3}\pi\times 27.5^3 = \qty{8.7e4}{nm^3}$।
\textbf{16.} \qty{60}{\%}: यानी बेलनों के सबसे कसे संभव जमाव का
दो-तिहाई --- \omterm{def:b1:nucleic-acids:chain}{DNA} लगभग क्रिस्टलीय कुंडलों में लिपटा है।
\textbf{17.} प्रति-आयनों (पॉलिऐमीन, $\mathrm{Mg^{2+}}$) को अधिकांश आवेश
उदासीन करना ही पड़ता है वरना प्रतिकर्षण पर पार न पाया जा सकता;
बचा प्रतिकर्षण और कसे कुंडलों में संचित मोड़ने की ऊर्जा
कवच पर बाहर की ओर ठेलते हैं, और जब पूँछ खुलती है तो यही दाब जीनोम का
पहला भाग बैक्टीरियम में ठेल देता है।
\textbf{18.} \qty{55}{nm} के किसी कवच में फंदों की त्रिज्याएँ
\qtyrange{10}{25}{nm} हैं, जो उस \qty{50}{nm} से नीचे है जिस पर \omterm{def:b1:nucleic-acids:chain}{DNA}
स्वच्छंद मुड़ता है: उसे मोड़ने में ऊर्जा पड़ती है, जो बाँधने वाला प्रेरक
(\omterm{def:b1:water-small-molecules:atp}{ATP}) चुकाता है और जो प्रश्न 17 के दाब के रूप में संचित हो जाती है।
\textbf{19.} $5'$-AGGTCGCCGCCC-$3'$।
\textbf{20.} \num{1200} \omterm{def:b1:nucleic-acids:nucleotide}{न्यूक्लियोटाइड}; $1200\times 330 = \qty{4.0e5}{Da}$।
\textbf{21.} $75\times 600 = \num{45000}$ युग्म, \qty{93}{\%} कूटलेखी
--- यानी एक सघन जीनोम; मानव जीनोम जीनों में साठ गुना कम सघन
है।
\textbf{22.} दो द्विशाखिकाएँ, हर एक \qty{1200}{s} में \num{24251} युग्म
की नकल करती: प्रति द्विशाखिका प्रति सेकंड \qty{20}{} युग्म (पोषी की
पॉलिमरेज़ \qty{1000}{} कर सकती है; समय समूचे चक्र से तय होता है,
पॉलिमरेज़ से नहीं)।
\textbf{23.} प्रति युग्म $\num{121207}/\num{48502} = 2.5$ बंध; प्रति सेकंड 20 युग्म पर दो द्विशाखिकाएँ: प्रति सेकंड 100 \omterm{def:b1:water-small-molecules:water}{हाइड्रोजन बंध} --- और उतने ही उनके पीछे फिर से बनते हुए।
\textbf{24.} \num{48502} में एक युग्म $G + C$ को
\qty{0.002}{\%} और $T_m$ को \qty{0.001}{\celsius} बदलता है: यानी समूचा
अणु पिघलाकर उसका पता ही नहीं चलता; केवल उस स्थल पर फैली कोई छोटी जाँच ही
उसे दिखाती।
\textbf{25.} \qty{0.1}{mol/L} सोडियम में $T_m = \qty{85}{\celsius}$;
परिरेखीय लंबाई \qty{16.5}{\micro m}, यानी कवच के व्यास की तीन सौ गुनी,
और उसी में उसके आयतन के \qty{60}{\%} पर बँधी हुई।
\end{solution}
